\section{The compact Hall factorization object at the Igusa boundary}
\label{sec:compact-hall-factorization-object}

\providecommand{\Z}{\mathbb Z}
\providecommand{\GammaBPS}{\Gamma_{\mathrm{BPS}}}
\providecommand{\Gammaeff}{\Gamma_{\mathrm{eff}}}
\providecommand{\Hall}{\mathcal H}
\providecommand{\Prim}{\operatorname{Prim}}
\providecommand{\sdim}{\operatorname{sdim}}
\providecommand{\BM}{\operatorname{BM}}
\providecommand{\Crit}{\operatorname{Crit}}
\providecommand{\van}{\varphi}
\providecommand{\eps}{\varepsilon}

This section records exactly what can be constructed from the present
input, and exactly where the compact construction stops.  The claim
under attack is the strongest possible one: that the compact
\(K3\times E\) BPS Hall/factorization object itself has been constructed
and that its protected primitive bracket is the denominator Lie
superalgebra \(\mathfrak g_{\Delta_5}\).  The conclusion is sharper:
there is a canonical finite-first construction once a compact oriented
critical Hall atlas is supplied, and its protected determinant is forced
to be the Igusa denominator once the primitive protected multiplicities
are the coefficients of \(\phi_{0,1}\).  The atlas, the primitive
comparison, and the bracket comparison are not consequences of the
scalar OP/DT identity of Oberdieck--Pixton and
Oberdieck--Pandharipande \cite{OB}, \cite{OPand}, or of the Borcherds
denominator theorem \cite{B}, \cite{GN}.

\subsection{The charge datum}

The Igusa boundary uses the charge lattice
\[
  \GammaBPS=\Z^3,\qquad
  \gamma=(n,l,m),
\]
with effective cone
\[
  \Gammaeff=
  \{m>0,\ n\geq 0,\ l\in\Z\}
  \cup
  \{m=0,\ n>0,\ l\in\Z\}
  \cup
  \{m=n=0,\ l<0\}.
\]
The embedding into the Borcherds lattice is
\[
  \alpha(n,l,m)=2n f_2-l f_3+2m f_{-2},
\]
and the quadratic convention is
\[
  (\alpha(\gamma),\alpha(\gamma))=-2(4nm-l^2).
\]
The protected signed multiplicity prescribed by the Borcherds product is
\[
  I_{\Delta_5}(n,l,m)=f(nm,l),
\]
where
\[
  \phi_{0,1}(\tau,z)=\sum_{n,l} f(n,l)q^n r^l .
\]
With the Weyl vector
\[
  \rho=f_2-\frac12 f_3+f_{-2},
\]
the denominator normalization used in the manuscript is
\[
  \operatorname{den}(\mathfrak g_{\Delta_5})
  =
  e^{-2\pi i(\rho,z)}
  \prod_{\gamma\in\Gammaeff}
  \left(1-e^{-2\pi i(\alpha(\gamma),z)}\right)^{f(nm,l)}
  =
  \frac{1}{64}\Delta_5(2Z).
\]
Equivalently, in the scalar variables \(q,r,s\),
\[
  \Delta_5(Z)
  =
  64(qrs)^{1/2}
  \prod_{\gamma=(n,l,m)\in\Gammaeff}
  (1-q^n r^l s^m)^{f(nm,l)} .
\]
These formulas are the denominator input of Borcherds and
Gritsenko--Nikulin \cite{B}, \cite{GN}.  They are not a construction of
the compact Hall object.

\subsection{The missing compact Hall atlas}

A compact oriented critical Hall atlas for \(X=K3\times E\), in a
stability sector \(\sigma\) and with orientation \(o\), consists of the
following finite-first data.

\begin{enumerate}
\item A cofinal system of finite lower charge sets
  \[
    \Gamma_{N,R}\subset\Gammaeff,\qquad
    \bigcup_{N,R}\Gamma_{N,R}=\Gammaeff,
  \]
  closed under passage to charges which occur in Harder--Narasimhan
  factors inside the sector.  Equivalently, if
  \(\gamma+\eta\in\Gamma_{N,R}\), then
  \(\gamma,\eta\in\Gamma_{N,R}\).  Thus the complement of
  \(\Gamma_{N,R}\) is a Hall ideal and quotient truncation is defined.
  The parameters \(N,R\) are auxiliary cutoffs; no invariant statement
  may depend on their particular definition.
\item For each \(\gamma\in\Gamma_{N,R}\), a derived moduli stack
  \(\mathfrak M_{\sigma,\gamma}\) of compact BPS objects, equipped
  locally on a Dwyer--Weiss--Williams/Cech/Ran cover by presentations
  \[
    \mathfrak M_{\sigma,\gamma}|_U\simeq
    [\Crit(f_{\gamma,U})/G_{\gamma,U}].
  \]
\item An orientation line \(\mathscr L_{o,\gamma}\) and compatible
  square-root data for the virtual canonical lines, transported through
  Thom--Sebastiani isomorphisms.
\item Extension correspondences
  \[
    \mathfrak M_{\sigma,\gamma}\times\mathfrak M_{\sigma,\eta}
    \xleftarrow{\ p\ }
    \mathfrak E_{\gamma,\eta}
    \xrightarrow{\ q\ }
    \mathfrak M_{\sigma,\gamma+\eta}
  \]
  with the compact-support/properness and Beck--Chevalley hypotheses
  needed for Borel--Moore pull-push on vanishing cycles.
\item Cech/Ran descent and Harder--Narasimhan completion data making the
  local Hall products independent of the chosen critical chart.
\end{enumerate}

This is precisely the compact input.  It is not supplied by the OP
scalar partition function, and it is not supplied by the formal
Borcherds algebra.

In the Vol.~III positive-cone language, this atlas is the Igusa-boundary
specialization of the finite-first oriented critical Hall cosheaf.  The
decategorification functor there can be used only after the compact
Hall cosheaf, orientation, and descent maps have been built; it is not a
substitute for them.

\begin{definition}[Finite compact Hall factorization object]
\label{def:finite-compact-hall-object}
Assume a compact oriented critical Hall atlas has been given on the
finite charge set \(\Gamma_{N,R}\).  Define
\[
  \mathsf F_{N,R}\Hall^{\mathrm{or}}_{X,\sigma,o}
  =
  \bigoplus_{\gamma\in\Gamma_{N,R}}
  H^{\BM}_{G_{\gamma}}
  \bigl(\Crit(f_{\gamma}),\van_{f_\gamma}\otimes
  \mathscr L_{o,\gamma}\bigr)[s_\gamma](t_\gamma),
\]
where the shift \(s_\gamma\) and Tate twist \(t_\gamma\) are the
normalizing ones required by the chosen orientation convention.  This is
the quotient truncation by the Hall ideal of charges outside
\(\Gamma_{N,R}\).  For
\(\gamma,\eta,\gamma+\eta\in\Gamma_{N,R}\), multiplication is the
composition
\[
  p^*(-)\ \xrightarrow{\mathrm{TS}}\ 
  \text{vanishing cycles on }\mathfrak E_{\gamma,\eta}
  \ \xrightarrow{\ q_!\ }\ 
  H^{\BM}_{G_{\gamma+\eta}}
  \bigl(\Crit(f_{\gamma+\eta}),\van_{f_{\gamma+\eta}}\otimes
  \mathscr L_{o,\gamma+\eta}\bigr),
\]
with the orientation transport inserted in the Thom--Sebastiani arrow.
If \(\gamma+\eta\notin\Gamma_{N,R}\), the product is zero in the
finite quotient and is computed in any larger cutoff containing
\(\gamma+\eta\).
\end{definition}

\begin{proposition}[What the atlas constructs]
\label{prop:compact-hall-atlas-constructs}
Under the hypotheses in Definition~\ref{def:finite-compact-hall-object},
\(\mathsf F_{N,R}\Hall^{\mathrm{or}}_{X,\sigma,o}\) is a finite
critical Hall truncation in the sense of Kontsevich--Soibelman
\cite{KSCoHA} and Davison--Meinhardt \cite{DavisonMeinhardtCoHA}.  The
Hall products are associative up to the
coherent homotopies supplied by the two-step extension stack.  If the
atlas satisfies Cech/Ran descent, the finite objects assemble into a
sectorial factorization cosheaf, and the completed compact Hall object is
\[
  \widehat{\Hall}^{\mathrm{or}}_{X,\sigma,o}
  =
  \varprojlim_{N,R}
  \mathsf F_{N,R}\Hall^{\mathrm{or}}_{X,\sigma,o}.
\]
\end{proposition}

\begin{proof}
The construction is formal once the atlas is present.  The product is
the usual Hall pull-push for extension correspondences, with
Thom--Sebastiani identifying the vanishing-cycle sheaf of a direct sum
with the external product of the two vanishing-cycle sheaves.  The
lower-set condition makes the discarded charges a Hall ideal, so the
quotient truncation inherits the product.  The
associativity constraint is the comparison between the two projections
from the stack of two-step flags.  The orientation line supplies the
sign and square-root correction.  Descent is exactly the assertion that
the local critical presentations and their extension correspondences
glue through the Cech/Ran nerve.  No Igusa-specific identity is used in
this proposition.
\end{proof}

\subsection{The determinant shadow}

The compact Hall object must not be identified with the denominator
algebra merely because its decategorified determinant has the same
product.  The determinant sees only signed primitive superdimensions.

\begin{definition}[Igusa protected integration]
\label{def:igusa-protected-integration}
An Igusa protected integration on
\(\widehat{\Hall}^{\mathrm{or}}_{X,\sigma,o}\) is a charge-preserving map
\[
  \chi^{\mathrm{prot}}_{\sigma,o}\colon
  K_0\!\left(\widehat{\Hall}^{\mathrm{or}}_{X,\sigma,o}\right)
  \longrightarrow
  \widehat{\mathbb T}_{\Gammaeff}
\]
to the completed charge torus, together with a primitive projection
\(\Prim_{\mathrm{prot}}\), such that
\[
  \sdim\,
  \Prim_{\mathrm{prot}}
  \left(\widehat{\Hall}^{\mathrm{or}}_{X,\sigma,o}\right)_{(n,l,m)}
  =
  f(nm,l)
  \qquad
  ((n,l,m)\in\Gammaeff).
\]
The orientation character of this integration is required to be the
Borcherds character \(\nu_{\Delta_5}\).
\end{definition}

\begin{proposition}[The exact determinant consequence]
\label{prop:compact-hall-determinant-shadow}
Assume the compact Hall object of
Proposition~\ref{prop:compact-hall-atlas-constructs} and an Igusa
protected integration in the sense of
Definition~\ref{def:igusa-protected-integration}.  Then the protected
primitive Fock determinant is the Igusa product:
\[
  64(qrs)^{1/2}
  \prod_{\gamma=(n,l,m)\in\Gammaeff}
  (1-q^n r^l s^m)^{
  \sdim\Prim_{\mathrm{prot}}(\widehat{\Hall}^{\mathrm{or}}_{X,\sigma,o})_\gamma}
  =
  \Delta_5(Z).
\]
In the Borcherds lattice normalization, the same statement is
\[
  \operatorname{den}_{\mathrm{prot}}
  \left(\widehat{\Hall}^{\mathrm{or}}_{X,\sigma,o}\right)
  =
  \frac1{64}\Delta_5(2Z).
\]
\end{proposition}

\begin{proof}
By Definition~\ref{def:igusa-protected-integration}, the exponent at
\((n,l,m)\) is \(f(nm,l)\).  Substitution into the Borcherds product for
\(\Delta_5\) gives the first formula.  Replacing \((q,r,s)\) by the
exponentials of the lattice coordinates for \(\alpha(n,l,m)\) and
including the Weyl factor gives the denominator normalization
\(64^{-1}\Delta_5(2Z)\).
\end{proof}

\begin{corollary}[What the determinant does not prove]
\label{cor:determinant-not-bracket}
The hypotheses of Proposition~\ref{prop:compact-hall-determinant-shadow}
do not identify the protected Hall primitive bracket with
\(\mathfrak g_{\Delta_5}\).  They determine only the signed primitive
superdimension in each charge.
\end{corollary}

\begin{proof}
A denominator product records the integers
\(\sdim\Prim_{\mathrm{prot}}(-)_\gamma\).  It does not record the
even/odd dimensions separately, the structure constants of the primitive
commutator, the real-root Serre relations, the imaginary-root
orthogonality relations, or the radical quotient required in a
Borcherds-Kac-Moody presentation.  These are bracket data, not
determinant data.
\end{proof}

\subsection{The bracket-level realization criterion}

The compact object realizes the Igusa denominator algebra only after the
following extra comparison data have been constructed.

\begin{definition}[Igusa Hall--Borcherds comparison datum]
\label{def:igusa-hall-borcherds-comparison}
An Igusa Hall--Borcherds comparison datum consists of:
\begin{enumerate}
\item the compact Hall object
  \(\widehat{\Hall}^{\mathrm{or}}_{K3\times E,\sigma,o}\);
\item an Igusa protected integration as in
  Definition~\ref{def:igusa-protected-integration};
\item a charge-preserving primitive commutator
  \[
    [-,-]_{\Hall}\colon
    \Prim_{\mathrm{prot},\gamma}\otimes
    \Prim_{\mathrm{prot},\eta}
    \longrightarrow
    \Prim_{\mathrm{prot},\gamma+\eta};
  \]
\item distinguished primitive classes for the three real simple roots
  \(\delta_1,\delta_2,\delta_3\);
\item a proof that the Hall primitive commutator satisfies the
  Gritsenko--Nikulin/Borcherds defining relations for
  \(\mathfrak g_{\Delta_5}\), including the real Serre relations,
  imaginary simple-root relations, parity convention, and Weyl-vector
  normalization;
\item a continuous Hopf pairing, negative half, Cartan part, and
  radical quotient producing the completed Hall--Drinfeld double;
\item a comparison with the formal chiral current envelope
  \(U_E^{\mathrm{ch}}(\mathfrak g_{\Delta_5}\otimes\mathcal O_E)\).
\end{enumerate}
\end{definition}

\begin{theorem}[Conditional compact Hall realization]
\label{thm:conditional-compact-hall-realization}
If an Igusa Hall--Borcherds comparison datum exists, then the protected
primitive Lie superalgebra of
\(\widehat{\Hall}^{\mathrm{or}}_{K3\times E,\sigma,o}\) is
charge-preservingly isomorphic to the denominator algebra
\(\mathfrak g_{\Delta_5}\).  Its protected denominator is
\(64^{-1}\Delta_5(2Z)\), and its chiral current envelope is the formal
factorization algebra
\[
  U_E^{\mathrm{ch}}(\mathfrak g_{\Delta_5}\otimes\mathcal O_E).
\]
\end{theorem}

\begin{proof}
The comparison datum gives a primitive Hall Lie superalgebra graded by
\(\Gammaeff\), with signed root multiplicities \(f(nm,l)\), real simple
roots \(\delta_i\), and exactly the Borcherds defining relations for
\(\mathfrak g_{\Delta_5}\).  The defining presentation of the
Gritsenko--Nikulin automorphic correction therefore identifies the two
charge-graded Lie superalgebras.  The denominator formula follows from
Proposition~\ref{prop:compact-hall-determinant-shadow}.  Applying the
Beilinson--Drinfeld chiral enveloping construction \cite{BD} to the
resulting constant Lie algebra over \(E\) gives the displayed current
envelope.
\end{proof}

\subsection{The obstruction}

\begin{theorem}[Present obstruction]
\label{thm:compact-hall-present-obstruction}
The sources used in the manuscript do not, by themselves, construct the
compact Hall/factorization object for \(K3\times E\) whose protected
primitive bracket is \(\mathfrak g_{\Delta_5}\).
\end{theorem}

\begin{proof}
Each available theorem supplies only part of the required structure.
The Borcherds and Gritsenko--Nikulin results \cite{B}, \cite{GN}
construct the denominator Lie superalgebra and its automorphic product;
they do not construct a compact BPS moduli problem, a Hall convolution,
or a primitive BPS bracket.  The OP/DT formula \cite{OB}, \cite{OPand}
supplies a scalar partition function and its normalization; it does not
construct the oriented critical Hall atlas of
Definition~\ref{def:finite-compact-hall-object}, the primitive
projection of Definition~\ref{def:igusa-protected-integration}, or the
Hall commutator relations in
Definition~\ref{def:igusa-hall-borcherds-comparison}.  The
Kontsevich--Soibelman and Davison--Meinhardt constructions
\cite{KSCoHA}, \cite{DavisonMeinhardtCoHA} supply the local
critical-Hall technology once the required critical charts, orientation,
and extension correspondences are present; they do not identify the
compact \(K3\times E\) sector with the Igusa Borcherds algebra.  The
Beilinson--Drinfeld current envelope \cite{BD} constructs the formal
chiral algebra attached to a known Lie algebra; it does not construct
the compact microscopic Lie algebra.

Thus the missing assertions are exactly:
\begin{enumerate}
\item a global compact oriented critical Hall atlas for the relevant
  \(K3\times E\) BPS moduli problem;
\item Harder--Narasimhan and Cech/Ran descent for its sectorial
  completion;
\item an orientation character equal to \(\nu_{\Delta_5}\);
\item a protected primitive integration whose charge-\((n,l,m)\) signed
  dimension is \(f(nm,l)\);
\item a Hall-commutator proof of the Gritsenko--Nikulin/Borcherds
  relations;
\item a continuous Hopf pairing and radical quotient giving the
  Hall--Drinfeld double;
\item a global oriented comparison from the holomorphic Chern--Simons
  factorization theory to the compact Hall construction.
\end{enumerate}
Without these data, the honest conclusion is the denominator shadow, not
the microscopic bracket-level realization.
\end{proof}

\subsection{Monotone path to closure}

The construction can be closed without weakening any claim by proving
the following statements in order.

\begin{enumerate}
\item Construct the compact oriented critical Hall atlas of
  Definition~\ref{def:finite-compact-hall-object} for each finite
  charge set \(\Gamma_{N,R}\).
\item Prove Cech/Ran and Harder--Narasimhan descent, so that the finite
  objects assemble into
  \(\widehat{\Hall}^{\mathrm{or}}_{K3\times E,\sigma,o}\).
\item Compute the orientation character and prove that it is
  \(\nu_{\Delta_5}\).
\item Construct the protected primitive integration and prove, charge by
  charge, that
  \[
    \sdim\Prim_{\mathrm{prot},(n,l,m)}=f(nm,l).
  \]
\item Identify the three real simple Hall primitives and prove the full
  Borcherds relation set in the Hall commutator.
\item Construct the completed Hopf pairing, negative half, Cartan part,
  and radical quotient, then identify the resulting double with
  \(\mathfrak g_{\Delta_5}\).
\item Apply the Beilinson--Drinfeld chiral enveloping functor to obtain
  the compactly sourced factorization algebra over \(E\).
\end{enumerate}

Items \(1\)--\(3\) produce the compact Hall/factorization object.  Item
\(4\) produces the Igusa determinant shadow.  Items \(5\)--\(6\) produce
the bracket-level denominator algebra.  Item \(7\) produces the chiral
current envelope from the microscopic source rather than by formal
postulate.
